\setcounter{Exercise}{0}

\chapter*{Feuille d'exercices : Les boucles}

\begin{Exercise}[title={Tableau d'exécution}]

    \Question Donner ligne par ligne dans un tableau la valeur des variables.
    \begin{lstlisting}[style=pythonstyle, language=Python]
total = 100
for i in range(1, 5):
    total = 100 - i
    print(total)\end{lstlisting}
    \Question À quel calcul correspond ce programme ?
    \Question Réécrire le programme avec une boucle while.
    \Question Donner ligne par ligne dans un tableau la valeur des variables.
    \begin{lstlisting}[style=pythonstyle, language=Python]
total = 0
for i in range(1, 3):
    if i % 2 == 0:
        total = total * i
    else:
        total = total * -i
    print(total)\end{lstlisting}
    \Question À quel calcul correspond ce programme ?
    \Question Réécrire le programme avec une boucle while.


    \Question Donner ligne par ligne dans un tableau la valeur des variables.
    \begin{lstlisting}[style=pythonstyle, language=Python]
total = 0
for i in range(1, 3):
    for j in range(1, 4):
        ajout = i**j
        total = total + ajout
print("Total :", total)\end{lstlisting}
    \Question À quel calcul correspond ce programme ?
\end{Exercise}

% \begin{Exercise}[title={Étoile}]
% \Question Écrire un programme permettant d'afficher un carré de 5 étoiles sur 5 comme ci-dessous :
% \begin{center}
% *****\\
% *****\\
% *****\\
% *****\\
% *****
% \end{center}
% \Question Écrire un programme permettant d'afficher un carré de 5000 étoiles sur 5000.
% \Question Écrire un programme permettant d'afficher un sapin de noël comme ci-dessous :
% \begin{center}
%     *\\
%     *\\
%    ***\\
%   *****\\
%  *******\\
%    ***\\
%    ***\\
% \end{center}
% \end{Exercise}

\newpage



\begin{Exercise}[title={Sommes}]
\Question Écrire un programme en Python pour calculer : $1+2+3+....+100$.
\Question Écrire un programme en Python pour calculer : $1+3+5+....+99$.
\Question Écrire un programme en Python pour calculer : $-1+2-3+4-5+....-99+100$.
\end{Exercise}